A la cambra acceleradora de la figura, de 30,0~cm de llargària, els electrons entren per l'esquerra i surten per la dreta. Mentre estan dins la cambra es mouen amb un MRUA (moviment rectilini uniformement accelerat), amb una acceleració cap a la dreta de $1,20\times 10^{13}\ \mathrm{m\,s^{-2}}$. En aquesta situació, es poden negligir les forces gravitatòries i els efectes relativistes.

\figura[0.5]{fig1.png}

\begin{apartat}{1}
Calculeu el camp elèctric a l'interior de la cambra acceleradora. Indiqueu-ne també la direcció i el sentit.
\end{apartat}

\begin{apartat}{1}
Quina diferència de potencial hi ha entre les parets esquerra i dreta de la cambra? Quina està a un potencial més alt? Quanta energia guanya cada electró que travessa la cambra?
\end{apartat}

\dades{%
  $Q_\text{electró} = 1,60\times 10^{-19}\ \mathrm{C}$\\
  $m_\text{electró} = 9,11\times 10^{-31}\ \mathrm{kg}$}
